% ============================================================
%  preamble.tex — packages, macros, environments
%  Differential Equations course. See CLAUDE.md for the rules.
% ============================================================

% ---------- Fonts and encoding ----------
\usepackage[T1]{fontenc}
\usepackage{lmodern}
\usepackage[utf8]{inputenc}

% ---------- Page layout ----------
\usepackage[a4paper, margin=1in, headheight=14pt]{geometry}
\usepackage{microtype}

% ---------- Mathematics ----------
\usepackage{amsmath, amssymb, amsthm}
\usepackage{mathtools}
% NOTE: the `physics` package is deliberately not loaded. Its \dv/\pdv would
% compete with \der/\pder below, and it globally redefines \div, \Re, \Im.

% ---------- Lists, tables ----------
\usepackage{enumitem}
\usepackage{booktabs}

% ---------- Graphics ----------
\usepackage{tikz}
\usetikzlibrary{arrows.meta, calc, decorations.markings, positioning}
\usepackage{pgfplots}
\pgfplotsset{compat=1.18}
\usepgfplotslibrary{fillbetween}
\usepackage[most]{tcolorbox}

% ---------- Bibliography ----------
\usepackage[backend=biber, style=alphabetic, maxbibnames=99]{biblatex}
\addbibresource{refs.bib}

% ---------- Hyperlinks (late), cleveref (after hyperref) ----------
\usepackage[hidelinks, bookmarksnumbered]{hyperref}
\usepackage[capitalise, nameinlink]{cleveref}

% ============================================================
%  Notation macros (Lagrange y', y'' and Newton \dot{x}, \ddot{x}
%  need no macro).
% ============================================================
\newcommand{\der}[2]{\frac{\mathrm{d}#1}{\mathrm{d}#2}}
\newcommand{\dern}[3]{\frac{\mathrm{d}^{#3}#1}{\mathrm{d}#2^{#3}}}
\newcommand{\pder}[2]{\frac{\partial #1}{\partial #2}}
\newcommand{\pdern}[3]{\frac{\partial^{#3} #1}{\partial #2^{#3}}}
\newcommand{\R}{\mathbb{R}}
\newcommand{\Lap}{\mathcal{L}}
\DeclareMathOperator{\W}{W}   % Wronskian

% Step justification in worked examples, used as  ... &= ... \why{reason} \\
\newcommand{\why}[1]{\qquad\text{\small\itshape(#1)}}

% One-line announcement of a notation switch inside a chapter.
\newcommand{\notation}[1]{%
  \par\smallskip\noindent
  {\small\sffamily\bfseries Notation\,$\rightarrow$\,}{\small\sffamily #1}%
  \par\smallskip}

% ============================================================
%  Theorem-like environments (numbered within chapter, shared counter)
% ============================================================
\theoremstyle{plain}
\newtheorem{theorem}{Theorem}[chapter]
\newtheorem{lemma}[theorem]{Lemma}
\newtheorem{corollary}[theorem]{Corollary}
\newtheorem{proposition}[theorem]{Proposition}

\theoremstyle{definition}
\newtheorem{definition}[theorem]{Definition}
\newtheorem{example}[theorem]{Example}

\theoremstyle{remark}
\newtheorem{remark}[theorem]{Remark}

% `proof` is provided by amsthm.

\crefname{theorem}{Theorem}{Theorems}
\crefname{lemma}{Lemma}{Lemmas}
\crefname{corollary}{Corollary}{Corollaries}
\crefname{proposition}{Proposition}{Propositions}
\crefname{definition}{Definition}{Definitions}
\crefname{example}{Example}{Examples}
\crefname{remark}{Remark}{Remarks}

% ============================================================
%  Boxed environments
% ============================================================
\definecolor{pitfallcol}{RGB}{176,48,48}
\definecolor{problemcol}{RGB}{40,80,140}
\definecolor{diagcol}{RGB}{90,90,90}

% --- pitfall: a common mistake and why it fails ---
\newcounter{pitfall}[chapter]
\renewcommand{\thepitfall}{\thechapter.\arabic{pitfall}}
\crefname{pitfall}{Pitfall}{Pitfalls}
\NewDocumentEnvironment{pitfall}{o}{%
  \refstepcounter{pitfall}%
  \begin{tcolorbox}[enhanced, breakable,
      colback=pitfallcol!4, colframe=pitfallcol!75!black,
      fonttitle=\bfseries\sffamily,
      title={Pitfall \thepitfall\IfValueT{#1}{: #1}}]%
}{\end{tcolorbox}}

% --- problem: tiered (W, C, H, P), numbered per tier within chapter ---
%   \begin{problem}{H}[Optional title or source]  ...  \end{problem}
\newcounter{probW}[chapter] \renewcommand{\theprobW}{\thechapter.W\arabic{probW}}
\newcounter{probC}[chapter] \renewcommand{\theprobC}{\thechapter.C\arabic{probC}}
\newcounter{probH}[chapter] \renewcommand{\theprobH}{\thechapter.H\arabic{probH}}
\newcounter{probP}[chapter] \renewcommand{\theprobP}{\thechapter.P\arabic{probP}}
\crefname{probW}{Problem}{Problems}
\crefname{probC}{Problem}{Problems}
\crefname{probH}{Problem}{Problems}
\crefname{probP}{Problem}{Problems}
\NewDocumentEnvironment{problem}{m o}{%
  \refstepcounter{prob#1}%
  \begin{tcolorbox}[enhanced, breakable,
      colback=white, colframe=problemcol,
      boxrule=0.6pt, arc=1mm,
      fonttitle=\bfseries\sffamily, coltitle=white,
      title={Problem \csname theprob#1\endcsname\IfValueT{#2}{\ \textmd{(#2)}}}]%
}{\end{tcolorbox}}

% --- diagnostic: the 3 prerequisite questions opening each chapter ---
\NewDocumentEnvironment{diagnostic}{}{%
  \begin{tcolorbox}[enhanced, breakable,
      colback=diagcol!5, colframe=diagcol,
      fonttitle=\bfseries\sffamily,
      title={Diagnostic (no answers given; if any of these stalls you, revisit the prerequisite first)}]%
  \begin{enumerate}[label=\textbf{D\arabic*.}, leftmargin=*]%
}{\end{enumerate}\end{tcolorbox}}

% --- hints: graded levels for H and P problems ---
%   \hintfor{prob:chNN-name}
%   \begin{hintlevels} \item idea \item setup \item first major step \end{hintlevels}
\newcommand{\hintfor}[1]{\par\medskip\noindent\textbf{Hints for \cref{#1}.}\par}
\newlist{hintlevels}{enumerate}{1}
\setlist[hintlevels]{label=\textsf{Level \arabic*.}, leftmargin=*, itemsep=2pt}

% ============================================================
%  pgfplots defaults for direction fields, phase lines, portraits
% ============================================================
\pgfplotsset{
  deplot/.style={
    axis lines=middle, grid=major, grid style={gray!20},
    xlabel style={anchor=west}, ylabel style={anchor=south},
    width=0.7\linewidth, height=0.5\linewidth,
    every axis plot/.append style={thick}},
  % Direction field: put in the axis options, then draw with
  %   \addplot3[slopearrows, quiver={u={1/sqrt(1+(F)^2)}, v={(F)/sqrt(1+(F)^2)},
  %     scale arrows=...}] (x,y,0);   where F is the right-hand side f(x,y).
  % Arrows have unit length before scaling, so steep slopes are not stretched.
  slopefield/.style={view={0}{90}, zmin=-1, zmax=1, samples=15},
  slopearrows/.style={gray, thin, -{Stealth[length=2pt]}},
}
